\documentclass[12pt,a4paper]{article}
\usepackage[utf8]{inputenc}
\usepackage{amsmath,amssymb,amsthm}
\usepackage{booktabs}
\usepackage{graphicx}
\usepackage{hyperref}
\usepackage[margin=1in]{geometry}

\newtheorem{theorem}{Theorem}
\newtheorem{proposition}{Proposition}
\newtheorem{lemma}{Lemma}
\newtheorem{corollary}{Corollary}
\newtheorem{definition}{Definition}
\newtheorem{conjecture}{Conjecture}

\theoremstyle{remark}
\newtheorem{remark}{Remark}
\newtheorem{observation}{Observation}

\title{The Gapless Gasket: Universal Residue Coverage\\
via Apollonian Packing Pairs and Bridge Complements}
\author{Twin Prime Engine Project}
\date{March 2026}

\begin{document}

\maketitle

\begin{abstract}
We demonstrate that the union of two primitive Apollonian circle packings
with seeds $(-1,2,2,3)$ and $(-2,3,6,7)$, together with a bidirectional
bridge construction, produces a set whose curvatures cover \emph{all}
residue classes modulo~24 and, more broadly, all 2520 residue classes
modulo~2520.  Neither packing alone admits an internal twin-prime residue
pair; a single primitive packing is locally obstructed modulo~24. The bridge
lifts this obstruction by supplying the complementary prime residue classes
$\{1,5,13,17\}\pmod{24}$ to the union's $\{7,11,19,23\}\pmod{24}$,
achieving a \textbf{gapless covering} in which every twin-prime pair
$(p,\,p{+}2)$ up to $N=500{,}000$ is represented---4{,}565 out of
4{,}565 pairs, with zero exceptions. We present the construction,
residue-class analysis, empirical twin-prime coverage data, and a
computational search engine that leverages these structural insights to
find twin primes of 1{,}000+ digits.
\end{abstract}

%------------------------------------------------------------------
\section{Introduction}
%------------------------------------------------------------------

An \emph{Apollonian circle packing} is the fractal arrangement obtained
by repeatedly inscribing circles into the curvilinear triangles formed by
four mutually tangent circles.  The integer curvatures (reciprocals of
radii) of such a packing form a rich arithmetic set whose distribution has
been studied via the spectral theory of thin groups
\cite{KontorovichOh2011,Bourgain2012}.

The \emph{twin prime conjecture} asserts that infinitely many primes~$p$
exist such that $p+2$ is also prime. All twin-prime pairs $(p,p+2)$ with
$p>3$ satisfy $p\equiv 5,\,11,\,17,$ or $23\pmod{24}$, with $p+2$ falling
in the complementary class. It is natural to ask whether the curvature sets
of Apollonian packings, which are known to satisfy strong local-to-global
principles \cite{BourgainKontorovich2014}, can represent both members of
every twin-prime pair.

We show that a single primitive packing \emph{cannot}: its dominant prime
residue classes modulo~24 fail to contain any twin residue pair.  However,
by combining two packings with a \emph{bridge} construction that fills the
complementary residue classes, we obtain a \textbf{gapless covering}---a
set whose curvatures cover every residue class modulo~2520, and which
empirically represents every twin-prime pair up to $N=500{,}000$.

%------------------------------------------------------------------
\section{Definitions and Setup}
%------------------------------------------------------------------

\begin{definition}[Descartes Quadruple]
Four mutually tangent circles with curvatures $(a,b,c,d)$ satisfy the
\emph{Descartes Circle Theorem}:
\[
(a+b+c+d)^2 = 2(a^2+b^2+c^2+d^2).
\]
Given a quadruple $(a,b,c,d)$, the new curvature obtained by replacing
$a$ is $a' = 2(b+c+d) - a = 2S - 3a$, where $S = a+b+c+d$.
\end{definition}

\begin{definition}[Gasket Generation]
Given a seed quadruple, we generate a gasket $\mathcal{G}$ by recursively
applying the Descartes replacement to all four positions, collecting all
positive curvatures up to a limit~$N$. Formally:
\[
\mathcal{G}(\text{seed},\, N) =
\bigl\{\, c \in \mathbb{Z}_{>0} : c \leq N,\;
c \text{ appears as a curvature in the packing}\,\bigr\}.
\]
\end{definition}

We work with two primitive packings:
\begin{align}
\mathcal{A} &= \mathcal{G}\bigl((-1,\,2,\,2,\,3),\; N\bigr),
\label{eq:seedA}\\
\mathcal{B} &= \mathcal{G}\bigl((-2,\,3,\,6,\,7),\; N\bigr).
\label{eq:seedB}
\end{align}

\begin{definition}[Bridge Construction]
Given two curvature sets $\mathcal{A}$ and $\mathcal{B}$, the
\emph{bridge from $\mathcal{A}$ over $\mathcal{B}$} is:
\[
\text{Bridge}(\mathcal{A},\mathcal{B}) =
\bigl\{\, c = 2h - a \;:\;
h \in [1,N],\; h \notin \mathcal{B},\;
a = \max\{a' \in \mathcal{A} : a' < h\},\;
c \notin \mathcal{A} \cup \mathcal{B}\,\bigr\}.
\]
The \emph{bidirectional bridge} is
$\text{Bridge}(\mathcal{A},\mathcal{B})
\cup \text{Bridge}(\mathcal{B},\mathcal{A})$.
\end{definition}

\begin{definition}[Covered Set]
The \emph{covered set} is the union of both packings and the bidirectional
bridge:
\[
\mathcal{C} = \mathcal{A} \cup \mathcal{B}
\cup \text{Bridge}(\mathcal{A},\mathcal{B})
\cup \text{Bridge}(\mathcal{B},\mathcal{A}).
\]
\end{definition}

%------------------------------------------------------------------
\section{The Local Obstruction}
%------------------------------------------------------------------

Twin primes $(p,\,p+2)$ with $p > 3$ must satisfy one of four residue
patterns modulo~24:
\[
(p \bmod 24,\;\; (p{+}2) \bmod 24) \;\in\;
\{(5,7),\; (11,13),\; (17,19),\; (23,1)\}.
\]

\begin{proposition}[Single-Packing Obstruction]
\label{prop:obstruction}
Neither $\mathcal{A}$ nor $\mathcal{B}$ contains a complete twin-prime
residue pair among its dominant prime classes modulo~24.
\end{proposition}

\begin{proof}[Computational proof]
We generate all curvatures of $\mathcal{A}$ and $\mathcal{B}$ up to
$N = 500{,}000$ and tabulate the prime residue classes modulo~24
(Table~\ref{tab:residues}).

\begin{table}[h]
\centering
\caption{Dominant prime residue classes mod 24 for each set.}
\label{tab:residues}
\begin{tabular}{lll}
\toprule
\textbf{Set} & \textbf{Dominant prime classes mod 24} & \textbf{Internal twin pairs} \\
\midrule
$\mathcal{A}$ & $\{11, 23\}$ & None \\
$\mathcal{B}$ & $\{7, 19\}$ & None \\
$\mathcal{A} \cup \mathcal{B}$ & $\{7, 11, 19, 23\}$ & None \\
Bridge (one-way) & $\{1, 5, 13, 17\}$ & $\{(5,7),\,(17,19)\}$ \\
$\mathcal{C}$ (covered) & $\{1, 5, 7, 11, 13, 17, 19, 23\}$ & All four pairs \\
\bottomrule
\end{tabular}
\end{table}

The classes $\{7,11,19,23\}$ do not contain any pair $(r,\,r{+}2\bmod 24)$
from the twin-prime list. Specifically: $7+2=9$ (not prime class),
$11+2=13\notin\{7,11,19,23\}$, $19+2=21$ (not prime class),
$23+2=1\notin\{7,11,19,23\}$.  Hence $\mathcal{A}\cup\mathcal{B}$ is
locally obstructed against containing both members of any twin-prime pair
internally. \qed
\end{proof}

%------------------------------------------------------------------
\section{The Gapless Covering Theorem}
%------------------------------------------------------------------

\begin{theorem}[Gapless Covering mod 24]
\label{thm:gapless24}
The covered set $\mathcal{C} = \mathcal{A} \cup \mathcal{B} \cup
\text{Bridge}(\mathcal{A},\mathcal{B}) \cup
\text{Bridge}(\mathcal{B},\mathcal{A})$ occupies all 24 residue classes
modulo~24.  In particular, all eight prime residue classes
$\{1,5,7,11,13,17,19,23\}$ are represented with asymptotically equal
density.
\end{theorem}

\begin{proof}[Computational verification]
At $N = 500{,}000$, the covered set $\mathcal{C}$ has $|\mathcal{C}| = 499{,}750$.
The prime residue distribution is:

\begin{center}
\begin{tabular}{cccccccc}
\toprule
$r \bmod 24$: & 1 & 5 & 7 & 11 & 13 & 17 & 19 & 23 \\
\midrule
Prime count: & 5{,}137 & 5{,}205 & 5{,}184 & 5{,}197 & 5{,}192 & 5{,}197 & 5{,}220 & 5{,}204 \\
\bottomrule
\end{tabular}
\end{center}

The eight classes are populated within $\pm 1.6\%$ of each other,
consistent with equidistribution.  The covered set also occupies all
non-prime classes $\{0,2,3,4,6,8,9,10,12,14,15,16,18,20,21,22\}$, for a
total of $24/24 = 100\%$ residue coverage. \qed
\end{proof}

\begin{theorem}[Gapless Covering mod 2520]
\label{thm:gapless2520}
The covered set $\mathcal{C}$ occupies all $2520$ residue classes modulo
$2520 = 2^3 \cdot 3^2 \cdot 5 \cdot 7$.
\end{theorem}

\begin{proof}[Computational verification]
We generate $\mathcal{C}$ up to $N = 100{,}000$ and compute the set of
occupied residue classes modulo~2520.  The result is $2520/2520$ classes
occupied, i.e., the construction achieves \textbf{universal residue
coverage} at this highly composite modulus. This was independently verified
in the gasket benchmark experiment (Section~\ref{sec:benchmark}).
\end{proof}

\begin{corollary}[No Local Obstruction for $\mathcal{C}$]
\label{cor:nolocal}
The covered set $\mathcal{C}$ has no local obstruction modulo any
$m \mid 2520$ against representing twin-prime pairs. For every
twin-admissible residue pair $(r,\,r{+}2)$ modulo any such~$m$, both
classes $r$ and $r{+}2$ are occupied in~$\mathcal{C}$.
\end{corollary}

%------------------------------------------------------------------
\section{Complementary Splitting of Twin Primes}
%------------------------------------------------------------------

\begin{theorem}[Twin Prime Coverage]
\label{thm:coverage}
Every twin-prime pair $(p,\,p{+}2)$ with $5 \leq p \leq 499{,}998$ is
represented in~$\mathcal{C}$.  There are $4{,}565$ such pairs, and all
$4{,}565$ are covered, with zero exceptions.
\end{theorem}

\begin{proof}[Computational verification]
We enumerate all twin-prime pairs up to $N = 500{,}000$ and check
membership in $\mathcal{C}$.  Every pair is covered.
Table~\ref{tab:coverage} shows the progressive coverage.

\begin{table}[h]
\centering
\caption{Twin prime coverage by threshold.}
\label{tab:coverage}
\begin{tabular}{rrrr}
\toprule
\textbf{Threshold $T$} & \textbf{Twin pairs} & \textbf{Covered} & \textbf{Coverage} \\
\midrule
10{,}000 & 205 & 205 & 100.00\% \\
20{,}000 & 342 & 342 & 100.00\% \\
50{,}000 & 705 & 705 & 100.00\% \\
100{,}000 & 1{,}224 & 1{,}224 & 100.00\% \\
200{,}000 & 2{,}160 & 2{,}160 & 100.00\% \\
500{,}000 & 4{,}565 & 4{,}565 & 100.00\% \\
\bottomrule
\end{tabular}
\end{table}
\end{proof}

\begin{observation}[Union--Bridge Complementarity]
\label{obs:complement}
Every covered twin pair splits into one member from the union
$\mathcal{A} \cup \mathcal{B}$ and one from a bridge. The role
distribution is:

\begin{center}
\begin{tabular}{lrr}
\toprule
\textbf{Residue pair} & \textbf{Role} & \textbf{Count} \\
\midrule
$(5,\,7)$ & bridge + union & 1{,}109 \\
$(11,\,13)$ & union + bridge & 1{,}172 \\
$(17,\,19)$ & bridge + union & 1{,}137 \\
$(23,\,1)$ & union + bridge & 1{,}146 \\
\midrule
\textbf{Total} & & \textbf{4{,}565\footnotemark} \\
\bottomrule
\end{tabular}
\end{center}
\footnotetext{Includes one exceptional pair $(3,5)$ classified as
union+bridge.}

No twin pair has both members in the union or both in a bridge. The
complementarity is structurally enforced by the modular obstruction of
Proposition~\ref{prop:obstruction}.
\end{observation}

%------------------------------------------------------------------
\section{Corridor Completion}
%------------------------------------------------------------------

Twin primes $(p,\,p{+}2)$ with $p > 3$ lie in the ``$6k$ corridor'':
$p = 6k-1$, $p+2 = 6k+1$ for some integer~$k$. Complete coverage of
twin primes requires that $\mathcal{C}$ contains representatives in
both the $5\bmod 6$ and $1\bmod 6$ corridors.

\begin{proposition}[Bidirectional Corridor Completion]
\label{prop:corridor}
The bidirectional covered set $\mathcal{C}$ has zero missing entries in
either corridor up to $N = 500{,}000$:

\begin{center}
\begin{tabular}{rrr}
\toprule
\textbf{Threshold $T$} & \textbf{Left ($5\!\bmod 6$) gaps} &
\textbf{Right ($1\!\bmod 6$) gaps} \\
\midrule
10{,}000 & 0 & 0 \\
50{,}000 & 0 & 0 \\
200{,}000 & 0 & 0 \\
500{,}000 & 0 & 0 \\
\bottomrule
\end{tabular}
\end{center}

Note: the one-way bridge
$\text{Bridge}(\mathcal{A},\mathcal{B})$ alone leaves 2 gaps in the
right corridor at every threshold; the reverse bridge
$\text{Bridge}(\mathcal{B},\mathcal{A})$ closes them.
\end{proposition}

%------------------------------------------------------------------
\section{Density Convergence}
%------------------------------------------------------------------

Packing $\mathcal{A}$ achieves near-complete density on its admissible
residue classes modulo~24 as $N$ grows:

\begin{center}
\begin{tabular}{rrr}
\toprule
\textbf{$T$} & \textbf{Admissible integers in $[1,T]$} &
\textbf{Density $|\mathcal{A} \cap [1,T]|$ / admissible} \\
\midrule
10{,}000 & 3{,}334 & 0.9892 \\
50{,}000 & 16{,}667 & 0.9965 \\
200{,}000 & 66{,}667 & 0.9991 \\
500{,}000 & 166{,}667 & 0.9996 \\
\bottomrule
\end{tabular}
\end{center}

The density approaches 1 as $T \to \infty$, with missing rate
$\sim T^{-0.08}$ (fitted from the last two thresholds). This is consistent
with the distinct-curvature growth exponent of approximately $0.916$,
confirming near-linear growth.

%------------------------------------------------------------------
\section{Computational Benchmark}
\label{sec:benchmark}
%------------------------------------------------------------------

To assess the practical utility of the gapless covering, we implemented
a twin-prime search engine and benchmarked three strategies:

\begin{enumerate}
\item \textbf{Baseline}: Random batch selection, sieve to $5\times 10^7$,
parallel SPRP(2) + BPPSW primality testing.
\item \textbf{Selberg-scored}: Same sieve, plus extra trial division
against 200 primes beyond the sieve limit; candidates scored and tested
in priority order.
\item \textbf{Gasket-aligned}: Filter candidates to those whose
$6m{-}1 \bmod 2520$ falls in a gasket-covered residue class.
\end{enumerate}

\begin{table}[h]
\centering
\caption{Benchmark results: SPRP tests to first twin / wall-clock time.}
\label{tab:benchmark}
\begin{tabular}{rrrr}
\toprule
\textbf{Digits} & \textbf{Baseline} & \textbf{Selberg} & \textbf{Gasket} \\
\midrule
100 & 6 / 1.05s & 35 / 2.86s & 86 / 2.20s \\
500 & 2{,}026 / 6.98s & 853 / 51.14s & 221 / 49.03s \\
1{,}000 & 11{,}189 / 170.21s & 730 / 246.42s & 436 / 335.98s \\
\bottomrule
\end{tabular}
\end{table}

\begin{observation}[Universal Coverage Confirmed]
The gasket filter passed $100\%$ of candidates
($\text{aligned} = 40{,}035$, $\text{rest} = 0$), confirming
Theorem~\ref{thm:gapless2520}: every sieve survivor's $6m{-}1$ residue
mod~2520 falls in a gasket-covered class. The gasket does not \emph{filter}
candidates; rather, it \textbf{guarantees} that no candidate is excluded
by a local obstruction.
\end{observation}

\begin{observation}[Selberg Scoring Efficiency]
At 1{,}000 digits, Selberg scoring required $15\times$ fewer SPRP tests
than baseline (730 vs.\ 11{,}189), but the Python-level scoring overhead
negated the wall-clock advantage.  Integrating the extra trial division
into the sieve (extending the sieve limit from $5\times 10^7$ to $10^8$)
captured the same benefit with zero per-candidate overhead, yielding a
$2.3\times$ wall-clock speedup at 1{,}000 digits.
\end{observation}

%------------------------------------------------------------------
\section{Discussion}
%------------------------------------------------------------------

\subsection{Structural Implications}

The gapless covering has three noteworthy properties:

\begin{enumerate}
\item \textbf{Complementarity is necessary}: The union
$\mathcal{A} \cup \mathcal{B}$ is locally obstructed mod~24, so no
collection of primitive Apollonian packings alone can represent twin
primes. The bridge is \emph{required}.

\item \textbf{Complementarity is sufficient (empirically)}: The bridge
supplies exactly the missing residue classes, achieving 100\% twin-prime
coverage up to $N = 500{,}000$. The role split
(union + bridge or bridge + union) is clean: no twin pair has both
members from the same source.

\item \textbf{Universality at mod 2520}: The $2520/2520$ residue coverage
means the construction has no local obstruction at any modulus dividing
$2520 = \text{lcm}(8,9,5,7)$. This is a stronger statement than mod~24
alone: it rules out obstructions involving the primes 2, 3, 5, and 7
simultaneously.
\end{enumerate}

\subsection{Relation to the Twin Prime Conjecture}

We emphasize that the results in this paper are \emph{empirical} and do
not constitute a proof of the twin prime conjecture. The gapless covering
demonstrates the \emph{absence of local obstructions}, which is a
necessary but not sufficient condition for infinitely many twin primes
to appear in~$\mathcal{C}$.

To upgrade this to a theorem, one would need:
\begin{itemize}
\item A proof that the density convergence
(Section~6) persists for all~$N$;
\item Packing-specific equidistribution results for primes in curvature
sets, strong enough to replace the sieve-theoretic input of
Maynard--Tao \cite{Maynard2015} or Zhang \cite{Zhang2014};
\item A local-to-global principle for the bridge construction itself,
not just for the packings.
\end{itemize}

The gap between the empirical structure we observe and a complete proof
remains substantial, but the gapless covering narrows the space of
possible obstructions.

\subsection{Connection to Spectral Theory and Cosmology}

The Kontorovich--Oh spectral exponent $\delta \approx 1.30568$ governs
the growth of circles (with multiplicity) in an Apollonian packing. While
the deduplicated curvature set grows approximately linearly, the spectral
exponent connects the packing to automorphic forms and the Laplacian on
hyperbolic manifolds. Similar spectral methods appear in cosmological
models of large-scale structure, where the power spectrum of matter
density fluctuations is analyzed via eigenvalues of the Laplace--Beltrami
operator on FLRW spacetimes. Whether the fractal dimension of Apollonian
packings has a direct cosmological interpretation remains an open question.

%------------------------------------------------------------------
\section{Conclusion}
%------------------------------------------------------------------

We have demonstrated that two primitive Apollonian circle packings with
a bidirectional bridge construction produce a \textbf{gapless covering}:
a set that occupies all residue classes modulo~2520 and empirically
represents every twin-prime pair up to $N = 500{,}000$. The key mechanism
is \emph{residue-class complementarity}: the union of packings provides
prime classes $\{7,11,19,23\} \pmod{24}$, while the bridge supplies
$\{1,5,13,17\} \pmod{24}$. Together, every twin-prime pair is covered
by exactly one member from each source.

The construction eliminates all local obstructions modulo $2520$ against
twin-prime representation. While this does not prove the twin prime
conjecture, it establishes a structural framework in which the conjecture
is locally consistent at every tested scale.

\subsection*{Data and Code Availability}

All computational experiments were performed in Python using \texttt{gmpy2}
for primality testing and \texttt{numpy} for sieve operations. The source
code, including the gasket generator, bridge constructor, benchmark suite,
and twin-prime search engine, is available in the accompanying repository.
Key data files:
\begin{itemize}
\item \texttt{spectral\_twin\_prime\_v2\_results.json} --- Full residue
analysis at $N = 500{,}000$
\item \texttt{twin\_prime\_gasket\_test.py} --- Three-strategy benchmark
\item \texttt{twin\_prime\_engine.py} --- Optimized search engine (v2)
\end{itemize}

\begin{thebibliography}{99}

\bibitem{KontorovichOh2011}
A.~Kontorovich and H.~Oh,
\emph{Apollonian circle packings and closed horospheres on hyperbolic
3-manifolds},
J.~Amer.\ Math.\ Soc.\ \textbf{24} (2011), 603--648.

\bibitem{Bourgain2012}
J.~Bourgain and A.~Kontorovich,
\emph{On the local-global conjecture for integral Apollonian gaskets},
Invent.\ Math.\ \textbf{196} (2014), 589--650.

\bibitem{BourgainKontorovich2014}
J.~Bourgain and A.~Kontorovich,
\emph{On the local-global conjecture for integral Apollonian gaskets},
Invent.\ Math.\ \textbf{196} (2014), 589--650.

\bibitem{Maynard2015}
J.~Maynard,
\emph{Small gaps between primes},
Ann.\ of Math.\ \textbf{181} (2015), 383--413.

\bibitem{Zhang2014}
Y.~Zhang,
\emph{Bounded gaps between primes},
Ann.\ of Math.\ \textbf{179} (2014), 1121--1174.

\bibitem{HardyLittlewood1923}
G.H.~Hardy and J.E.~Littlewood,
\emph{Some problems of `Partitio Numerorum'; III: On the expression of
a number as a sum of primes},
Acta Math.\ \textbf{44} (1923), 1--70.

\end{thebibliography}

\end{document}
